\documentclass[12pt,a4paper]{article}

\usepackage[a4paper,margin=1.7cm]{geometry}
\usepackage[T2A]{fontenc}
\usepackage[utf8]{inputenc}
\usepackage[russian]{babel}
\usepackage{amsmath,amssymb,mathtools}
\usepackage{array,booktabs,longtable}
\usepackage{xcolor}
\usepackage{hyperref}
\usepackage{tikz}

\definecolor{reviewteal}{RGB}{0,128,128}

\hypersetup{
    unicode=true,
    colorlinks=true,
    linkcolor=blue!55!black,
    urlcolor=blue!55!black
}

\setlength{\parindent}{0pt}
\setlength{\parskip}{0.55em}
\renewcommand{\arraystretch}{1.25}
\sloppy

\newcommand{\errorlabel}[1]{\textcolor{red}{Ошибка:} #1}
\newcommand{\keyidea}[1]{\textcolor{blue!60!black}{Ключевая идея:} #1}
\newcommand{\resultlabel}[1]{\textcolor{reviewteal!80!black}{Итог:} #1}

\begin{document}

\pagenumbering{gobble}

\begin{titlepage}
\thispagestyle{empty}
\centering

\vspace*{0.28\textheight}

{\LARGE\bfseries Ревью домашней работы\par}

\vspace{1.4cm}

{\large Автор: Лёвин Артём Александрович\par}

\vspace{0.7cm}

{\large 26 сентября 2026 года\par}

\vfill

\begin{tikzpicture}
\draw[blue!35!black,line width=0.7pt]
(0,0) .. controls (2,0.45) and (4,-0.45) .. (6,0);
\end{tikzpicture}

\end{titlepage}

\clearpage
\pagenumbering{arabic}

\section*{Сводная таблица}

\small

\begin{longtable}{
    >{\centering\arraybackslash}p{0.8cm}
    p{3.0cm}
    p{4.0cm}
    p{2.2cm}
    p{4.1cm}
}
\toprule
\textbf{№} &
\textbf{Тема} &
\textbf{Суть ошибки} &
\textbf{Ваш ответ} &
\textbf{Правильный ответ}
\\
\midrule
\endfirsthead

\toprule
\textbf{№} &
\textbf{Тема} &
\textbf{Суть ошибки} &
\textbf{Ваш ответ} &
\textbf{Правильный ответ}
\\
\midrule
\endhead

\hyperref[task:4]{4}
&
Векторы. Длина вектора
&
Найдены координаты результирующего вектора, но не вычислена его длина.
&
$(5;12)$
&
$13$
\\

\addlinespace

\hyperref[task:8]{8}
&
Условная вероятность
&
Вероятность положительного теста при наличии заболевания принята за вероятность заболевания при положительном результате.
&
$0{,}86$
&
$0{,}43$
\\

\addlinespace

\hyperref[task:15]{15}
&
Движение по течению и против течения
&
Уравнение составлено верно, но при умножении на общий знаменатель неверно учтено время движения.
&
$360$ км
&
$288$ км
\\

\addlinespace

\hyperref[task:19]{19}
&
Рациональное неравенство
&
После верного приведения к общему знаменателю неверно раскрыты скобки в числителе; решение не доведено до ответа.
&
не указан
&
\(\displaystyle
(-\infty;-4]\cup(-3;-1]\cup(0;7)\cup(7;+\infty)
\)
\\

\bottomrule
\end{longtable}

\normalsize

\newpage

\thispagestyle{empty}

\vspace*{\fill}

\begin{center}
{\LARGE\bfseries Разбор ошибок}
\end{center}

\vspace*{\fill}

% ============================================================
% ЗАДАНИЕ 4
% ============================================================

\newpage
\phantomsection
\label{task:4}

\section*{Задание 4}

\textbf{Текст задания.}

Даны векторы
\[
\vec a=(1;1),\qquad \vec b=(0;7).
\]

Найдите длину вектора
\[
5\vec a+\vec b.
\]

\textbf{Ваше решение.}

В работе получено:
\[
5\vec a+\vec b
=
5(1;1)+(0;7)
=
(5;5)+(0;7)
=
(5;12).
\]

\textbf{Ваш ответ.}

\[
(5;12).
\]

\errorlabel{ найдены координаты вектора, но не найдена его длина.}

\textbf{Где именно возникает ошибка.}

Последний правильный шаг:
\[
5\vec a+\vec b=(5;12).
\]

До этого момента решение выполнено верно. Координаты результирующего
вектора действительно равны $(5;12)$.

Ошибка возникает после этого: полученные координаты записываются как
окончательный ответ, хотя по условию требуется найти не сам вектор, а
его \emph{длину}.

Координаты вектора и длина вектора --- разные величины. Координаты
записываются парой чисел, а длина является одним неотрицательным числом.

\textbf{Почему этот шаг неверен.}

Если вектор имеет координаты
\[
\vec v=(x;y),
\]
то его длина находится по формуле
\[
|\vec v|=\sqrt{x^2+y^2}.
\]

Эта формула получается из теоремы Пифагора: горизонтальная и вертикальная
составляющие вектора образуют катеты прямоугольного треугольника, а сам
вектор соответствует его гипотенузе.

Поэтому после получения координат $(5;12)$ решение ещё не закончено.

\keyidea{ сначала необходимо найти координаты результирующего вектора, а затем отдельно вычислить его длину.}

\textbf{Исправление от последнего правильного шага.}

Мы уже знаем:
\[
5\vec a+\vec b=(5;12).
\]

Следовательно,
\[
|5\vec a+\vec b|
=
\sqrt{5^2+12^2}.
\]

Вычисляем:
\[
5^2=25,
\qquad
12^2=144,
\]
поэтому
\[
|5\vec a+\vec b|
=
\sqrt{25+144}
=
\sqrt{169}
=
13.
\]

\textbf{Полное правильное решение.}

Сначала умножаем вектор $\vec a$ на число $5$:
\[
5\vec a
=
5(1;1)
=
(5;5).
\]

Прибавляем вектор $\vec b$:
\[
5\vec a+\vec b
=
(5;5)+(0;7)
=
(5;12).
\]

Теперь находим длину полученного вектора:
\[
|5\vec a+\vec b|
=
\sqrt{5^2+12^2}
=
\sqrt{25+144}
=
\sqrt{169}
=
13.
\]

\textbf{Проверка.}

Проверим вычисление:
\[
5^2+12^2=25+144=169,
\]
а
\[
13^2=169.
\]

Следовательно, длина действительно равна $13$.

\resultlabel{ координаты вектора $(5;12)$ были найдены правильно, но требовалось выполнить ещё один шаг. Правильная длина вектора равна $13$.}

\textbf{Задача на закрепление.}

Даны векторы
\[
\vec a=(2;-1),
\qquad
\vec b=(-1;4).
\]
Найдите длину вектора
\[
3\vec a+\vec b.
\]

% ============================================================
% ЗАДАНИЕ 8
% ============================================================

\newpage
\phantomsection
\label{task:8}

\section*{Задание 8}

\textbf{Текст задания.}

При подозрении на наличие некоторого заболевания пациента отправляют
на ПЦР-тест. Если заболевание действительно есть, тест подтверждает его
в $86\%$ случаев. Если заболевания нет, тест выявляет его отсутствие
в среднем в $94\%$ случаев. Известно, что в среднем тест оказывается
положительным у $10\%$ пациентов, направленных на тестирование.

Тест некоторого пациента оказался положительным.

Какова вероятность того, что пациент действительно имеет это заболевание?

\textbf{Ваше решение.}

Промежуточные вычисления в предоставленном решении не записаны.
Указан только окончательный ответ.

\textbf{Ваш ответ.}

\[
0{,}86.
\]

\errorlabel{ перепутаны две разные условные вероятности.}

\textbf{Где именно возникает ошибка.}

В условии действительно присутствует число
\[
0{,}86.
\]

Но оно означает следующее:

если заранее известно, что человек болен, то вероятность положительного
результата теста равна $0{,}86$.

В задаче спрашивается обратное:

если уже известно, что тест положительный, какова вероятность того,
что человек болен?

Это два разных события с разным порядком условия. Они не имеют оснований
совпадать.

Именно поэтому непосредственно записывать $0{,}86$ в ответ нельзя.

\textbf{Почему это неверно.}

Положительный тест может появиться двумя способами:

\begin{enumerate}
    \item человек действительно болен, и тест правильно обнаружил заболевание;
    \item человек не болен, но тест дал ложноположительный результат.
\end{enumerate}

По условию у здорового пациента тест правильно показывает отсутствие
заболевания в $94\%$ случаев.

Значит, вероятность ложноположительного результата равна
\[
1-0{,}94=0{,}06.
\]

Таким образом, среди всех положительных результатов присутствуют как
действительно заболевшие, так и здоровые пациенты.

Нужно определить, какую часть всех положительных результатов составляют
истинно положительные результаты.

\keyidea{ вероятность положительного теста у заболевшего и вероятность заболевания у пациента с положительным тестом --- это разные величины. Сначала нужно восстановить долю заболевших среди всех тестируемых.}

\textbf{Исправление.}

Обозначим через $p$ долю заболевших среди всех направленных на тестирование.

Если человек болен, положительный результат получается с вероятностью
$0{,}86$.

Поэтому доля истинно положительных тестов равна
\[
0{,}86p.
\]

Если человек не болен, его доля составляет
\[
1-p.
\]

Ложноположительный результат среди здоровых возникает с вероятностью
\[
0{,}06.
\]

Поэтому доля ложноположительных результатов равна
\[
0{,}06(1-p).
\]

Всего положительный результат получают $10\%$ тестируемых, то есть
\[
0{,}10.
\]

Получаем уравнение:
\[
0{,}86p+0{,}06(1-p)=0{,}10.
\]

Раскрываем скобки:
\[
0{,}86p+0{,}06-0{,}06p=0{,}10.
\]

Приводим подобные слагаемые:
\[
0{,}80p+0{,}06=0{,}10.
\]

Следовательно,
\[
0{,}80p=0{,}04,
\]
откуда
\[
p=0{,}05.
\]

Значит, среди всех направленных на тестирование заболевание действительно
имеют $5\%$ пациентов.

Теперь найдём долю пациентов, которые одновременно болеют и получают
положительный результат:
\[
0{,}05\cdot0{,}86=0{,}043.
\]

При этом общая доля положительных тестов равна
\[
0{,}10.
\]

Значит, среди всех пациентов с положительным результатом доля действительно
заболевших равна
\[
\frac{0{,}043}{0{,}10}=0{,}43.
\]

\textbf{Полное правильное решение.}

Вероятность ложноположительного результата равна
\[
1-0{,}94=0{,}06.
\]

Пусть $p$ --- вероятность того, что случайно выбранный пациент из
направленных на тестирование действительно болен.

По формуле полной вероятности для положительного результата:
\[
0{,}86p+0{,}06(1-p)=0{,}10.
\]

Решаем:
\[
0{,}86p+0{,}06-0{,}06p=0{,}10,
\]
\[
0{,}80p=0{,}04,
\]
\[
p=0{,}05.
\]

Вероятность одновременно иметь заболевание и получить положительный
результат:
\[
0{,}05\cdot0{,}86=0{,}043.
\]

Вероятность положительного теста по условию равна
\[
0{,}10.
\]

Поэтому искомая вероятность:
\[
\frac{0{,}043}{0{,}10}=0{,}43.
\]

То есть
\[
43\%.
\]

\textbf{Проверка.}

Проверим, действительно ли при найденной доле заболевших положительный
результат возникает в $10\%$ случаев.

Истинно положительные результаты:
\[
0{,}05\cdot0{,}86=0{,}043.
\]

Здоровых пациентов:
\[
1-0{,}05=0{,}95.
\]

Ложноположительные результаты:
\[
0{,}95\cdot0{,}06=0{,}057.
\]

Всего:
\[
0{,}043+0{,}057=0{,}10.
\]

Получено ровно то значение, которое дано в условии.

\resultlabel{ вероятность того, что пациент с положительным тестом действительно имеет заболевание, равна $0{,}43$, то есть $43\%$.}

\textbf{Задача на закрепление.}

Если заболевание действительно есть, тест даёт положительный результат
в $80\%$ случаев. Если заболевания нет, тест даёт отрицательный результат
в $90\%$ случаев. В среднем положительный результат получают $24\%$
тестируемых. У некоторого пациента тест оказался положительным.
Найдите вероятность того, что пациент действительно имеет заболевание.

% ============================================================
% ЗАДАНИЕ 15
% ============================================================

\newpage
\phantomsection
\label{task:15}

\section*{Задание 15}

\textbf{Текст задания.}

Теплоход, скорость которого в неподвижной воде равна $15$ км/ч,
проходит по течению реки и после стоянки возвращается в исходный пункт.
Скорость течения равна $3$ км/ч, стоянка длится $5$ часов, а в исходный
пункт теплоход возвращается через $25$ часов после отплытия из него.

Сколько километров проходит теплоход за весь рейс?

\textbf{Ваше решение.}

Правильно найдены скорости теплохода.

По течению:
\[
15+3=18\text{ км/ч}.
\]

Против течения:
\[
15-3=12\text{ км/ч}.
\]

Далее весь путь обозначен через $n$. Тогда путь в одну сторону равен
половине всего пути.

Время движения по течению записано как
\[
\frac{n}{36},
\]
а время движения против течения --- как
\[
\frac{n}{24}.
\]

После этого составлено уравнение
\[
\frac{n}{36}+\frac{n}{24}+5=25.
\]

Этот шаг выполнен правильно.

Далее в работе появляется выражение, приводящее к
\[
5n=1800,
\]
после чего записывается
\[
n=360.
\]

\textbf{Ваш ответ.}

\[
360\text{ км}.
\]

\errorlabel{ при преобразовании верно составленного уравнения допущена арифметическая ошибка.}

\textbf{Где именно возникает ошибка.}

Последний полностью правильный шаг:
\[
\frac{n}{36}+\frac{n}{24}+5=25.
\]

Сначала нужно убрать из общего времени пятичасовую стоянку:
\[
\frac{n}{36}+\frac{n}{24}=25-5.
\]

То есть
\[
\frac{n}{36}+\frac{n}{24}=20.
\]

На фотографии далее при приведении к знаменателю $72$ фактически
используется число $25$, поскольку появляется число
\[
1800=25\cdot72.
\]

Но после переноса стоянки справа уже должно находиться число $20$.

Следовательно, при умножении на $72$ справа должно получиться
\[
20\cdot72=1440.
\]

\textbf{Почему этот шаг неверен.}

Полные $25$ часов состоят из двух частей:

\begin{itemize}
    \item времени движения;
    \item пяти часов стоянки.
\end{itemize}

Поэтому собственно на движение теплоход затратил
\[
25-5=20
\]
часов.

В составленном Вами уравнении это уже учтено правильно. Ошибка появляется
не в модели движения, а только при последующем арифметическом преобразовании.

Если после переноса числа $5$ продолжить использовать число $25$, то
пятичасовая стоянка фактически перестаёт вычитаться из времени движения.

\keyidea{ уравнение было составлено верно. Нужно сохранить равносильность каждого следующего преобразования: сначала получить двадцать часов движения, затем умножать обе части на общий знаменатель.}

\textbf{Исправление от последнего правильного шага.}

Начинаем с верного уравнения:
\[
\frac{n}{36}+\frac{n}{24}+5=25.
\]

Переносим $5$:
\[
\frac{n}{36}+\frac{n}{24}=20.
\]

Наименьший общий знаменатель чисел $36$ и $24$ равен $72$.

Умножаем всё уравнение на $72$:
\[
72\cdot\frac{n}{36}
+
72\cdot\frac{n}{24}
=
72\cdot20.
\]

Получаем:
\[
2n+3n=1440.
\]

Следовательно,
\[
5n=1440,
\]
откуда
\[
n=\frac{1440}{5}=288.
\]

\textbf{Полное правильное решение.}

Скорость теплохода по течению:
\[
15+3=18\text{ км/ч}.
\]

Скорость теплохода против течения:
\[
15-3=12\text{ км/ч}.
\]

Пусть $n$ километров --- весь путь теплохода за рейс.

Так как теплоход проходит одинаковое расстояние в одну сторону и обратно,
расстояние каждого участка равно
\[
\frac{n}{2}.
\]

Время движения по течению:
\[
\frac{\frac n2}{18}
=
\frac{n}{36}.
\]

Время движения против течения:
\[
\frac{\frac n2}{12}
=
\frac{n}{24}.
\]

Кроме движения была стоянка продолжительностью $5$ часов.

По условию от отплытия до возвращения прошло $25$ часов. Поэтому
\[
\frac{n}{36}+\frac{n}{24}+5=25.
\]

Вычитаем $5$:
\[
\frac{n}{36}+\frac{n}{24}=20.
\]

Умножаем на $72$:
\[
2n+3n=1440.
\]

Получаем:
\[
5n=1440,
\]
\[
n=288.
\]

\textbf{Проверка.}

Если весь путь равен $288$ км, то расстояние в каждую сторону равно
\[
\frac{288}{2}=144\text{ км}.
\]

Время движения по течению:
\[
\frac{144}{18}=8\text{ ч}.
\]

Время движения против течения:
\[
\frac{144}{12}=12\text{ ч}.
\]

Добавляем стоянку:
\[
8+5+12=25\text{ ч}.
\]

Получено ровно указанное в условии время.

\resultlabel{ верная математическая модель в решении уже была получена. Ошибка появилась только при арифметическом преобразовании. За весь рейс теплоход проходит $288$ км.}

\textbf{Задача на закрепление.}

Скорость катера в неподвижной воде равна $12$ км/ч, а скорость течения
реки --- $4$ км/ч. Катер идёт по течению, стоит $2$ часа и возвращается
обратно. От момента выхода до возвращения проходит $11$ часов.
Найдите весь путь катера за рейс.

% ============================================================
% ЗАДАНИЕ 19
% ============================================================

\newpage
\phantomsection
\label{task:19}

\section*{Задание 19}

\textbf{Текст задания.}

Решите неравенство
\[
\frac{3x+28-x^2}{x^2-7x}
\le
2+\frac{2}{x+3}.
\]

\textbf{Ваше решение.}

В работе правильно учтены ограничения:
\[
x\ne -3,
\qquad
x\ne0,
\qquad
x\ne7.
\]

Далее выражение приводится к общему знаменателю:
\[
\frac{
(3x+28-x^2)(x+3)
-2(x^2-7x)
-2(x^2-7x)(x+3)
}{
(x^2-7x)(x+3)
}
\le0.
\]

Этот этап выполнен правильно.

После раскрытия скобок в работе получается многочлен
\[
-3x^3-10x^2+9x+84,
\]
после чего начинается подбор его целых корней.

Окончательного решения неравенства и окончательного ответа в работе нет.

\textbf{Ваш ответ.}

не указан.

\errorlabel{ неверно раскрыты скобки и приведены подобные слагаемые в числителе после правильного приведения к общему знаменателю.}

\textbf{Где именно возникает ошибка.}

Последний правильный шаг:
\[
\frac{
(3x+28-x^2)(x+3)
-2(x^2-7x)
-2(x^2-7x)(x+3)
}{
(x^2-7x)(x+3)
}
\le0.
\]

После этого нужно аккуратно раскрыть каждую группу скобок отдельно.

Рассмотрим первое произведение:
\[
(3x+28-x^2)(x+3).
\]

Раскрываем:
\[
3x(x+3)+28(x+3)-x^2(x+3).
\]

Получаем:
\[
3x^2+9x+28x+84-x^3-3x^2.
\]

Слагаемые $3x^2$ и $-3x^2$ сокращаются:
\[
(3x+28-x^2)(x+3)
=
-x^3+37x+84.
\]

Вторая часть:
\[
-2(x^2-7x)
=
-2x^2+14x.
\]

Для третьей части сначала раскрываем:
\[
(x^2-7x)(x+3)
=
x^3+3x^2-7x^2-21x.
\]

То есть
\[
(x^2-7x)(x+3)
=
x^3-4x^2-21x.
\]

Умножаем на $-2$:
\[
-2(x^2-7x)(x+3)
=
-2x^3+8x^2+42x.
\]

Теперь складываем:
\[
(-x^3+37x+84)
+
(-2x^2+14x)
+
(-2x^3+8x^2+42x).
\]

Получаем:
\[
-3x^3+6x^2+93x+84.
\]

Следовательно, записанный в работе многочлен
\[
-3x^3-10x^2+9x+84
\]
не является числителем исходного выражения.

\textbf{Почему это неверно.}

В рациональном неравенстве вся дальнейшая работа методом интервалов
зависит от нулей числителя и знаменателя.

Если при раскрытии скобок получается другой многочлен, то меняются его
корни, меняются критические точки и, следовательно, меняется всё решение
неравенства.

Поэтому подбор корней после ошибочного раскрытия скобок уже относится
не к исходному неравенству.

\keyidea{ после приведения рационального неравенства к одной дроби нужно раскрывать крупные произведения по отдельности, затем приводить подобные слагаемые и только после этого раскладывать числитель и знаменатель на множители. При сокращении запрещённые значения исходного выражения необходимо сохранять.}

\textbf{Исправление от последнего правильного шага.}

После правильного приведения к общей дроби имеем:
\[
\frac{-3x^3+6x^2+93x+84}
{x(x-7)(x+3)}
\le0.
\]

Раскладываем числитель на множители.

Сначала выносим $-3$:
\[
-3x^3+6x^2+93x+84
=
-3(x^3-2x^2-31x-28).
\]

Проверим число $7$:
\[
7^3-2\cdot7^2-31\cdot7-28
=
343-98-217-28
=
0.
\]

Значит, $x-7$ является множителем.

После разложения получаем:
\[
x^3-2x^2-31x-28
=
(x-7)(x+1)(x+4).
\]

Следовательно,
\[
-3x^3+6x^2+93x+84
=
-3(x-7)(x+1)(x+4).
\]

Знаменатель:
\[
x^2-7x=x(x-7),
\]
поэтому
\[
(x^2-7x)(x+3)
=
x(x-7)(x+3).
\]

Получаем:
\[
\frac{-3(x-7)(x+1)(x+4)}
{x(x-7)(x+3)}
\le0.
\]

По области допустимых значений
\[
x\ne7.
\]

Поэтому множитель $x-7$ можно сократить:
\[
\frac{-3(x+1)(x+4)}
{x(x+3)}
\le0.
\]

Но значение $x=7$ после сокращения нельзя вернуть в область допустимых
значений. Оно по-прежнему запрещено исходным выражением.

Делим на отрицательное число $-3$. При этом знак неравенства меняется:
\[
\frac{(x+1)(x+4)}
{x(x+3)}
\ge0.
\]

Основные критические точки:
\[
-4,\qquad -3,\qquad -1,\qquad 0.
\]

Точки $-4$ и $-1$ обращают числитель в ноль.

Точки $-3$ и $0$ обращают знаменатель в ноль.

Кроме того, отдельно сохраняется запрещённое значение
\[
x=7.
\]

\textbf{Полное правильное решение.}

Исходное неравенство:
\[
\frac{3x+28-x^2}{x^2-7x}
\le
2+\frac{2}{x+3}.
\]

Сначала выписываем область допустимых значений:
\[
x^2-7x\ne0,
\qquad
x+3\ne0.
\]

Следовательно,
\[
x\ne0,\qquad x\ne7,\qquad x\ne-3.
\]

Переносим всё в левую часть:
\[
\frac{3x+28-x^2}{x^2-7x}
-
2
-
\frac{2}{x+3}
\le0.
\]

Приводим к общему знаменателю:
\[
\frac{
(3x+28-x^2)(x+3)
-2(x^2-7x)(x+3)
-2(x^2-7x)
}{
(x^2-7x)(x+3)
}
\le0.
\]

Раскрываем скобки:
\[
(3x+28-x^2)(x+3)
=
-x^3+37x+84,
\]
\[
-2(x^2-7x)(x+3)
=
-2x^3+8x^2+42x,
\]
\[
-2(x^2-7x)
=
-2x^2+14x.
\]

Поэтому числитель равен
\[
-3x^3+6x^2+93x+84.
\]

Получаем:
\[
\frac{-3x^3+6x^2+93x+84}
{x(x-7)(x+3)}
\le0.
\]

Раскладываем числитель:
\[
-3x^3+6x^2+93x+84
=
-3(x-7)(x+1)(x+4).
\]

Следовательно,
\[
\frac{-3(x-7)(x+1)(x+4)}
{x(x-7)(x+3)}
\le0.
\]

При $x\ne7$ сокращаем общий множитель:
\[
\frac{-3(x+1)(x+4)}
{x(x+3)}
\le0.
\]

После деления на $-3$:
\[
\frac{(x+1)(x+4)}
{x(x+3)}
\ge0.
\]

Располагаем критические точки по возрастанию:
\[
-4<-3<-1<0.
\]

Рассматриваем промежутки:
\[
(-\infty;-4),\quad
(-4;-3),\quad
(-3;-1),\quad
(-1;0),\quad
(0;+\infty).
\]

Знак выражения на них соответственно:
\[
+,\quad -,\quad +,\quad -,\quad +.
\]

Так как требуется значение не меньше нуля, берём промежутки с положительным
знаком и нули числителя:
\[
(-\infty;-4]
\cup
(-3;-1]
\cup
(0;+\infty).
\]

Но исходная область допустимых значений дополнительно запрещает $x=7$.
Поэтому последний промежуток разбивается на два:

\[
(0;7)\cup(7;+\infty).
\]

Окончательно:
\[
\boxed{
(-\infty;-4]
\cup
(-3;-1]
\cup
(0;7)
\cup
(7;+\infty)
}.
\]

\textbf{Проверка.}

Точки
\[
x=-4
\quad\text{и}\quad
x=-1
\]
обращают числитель в ноль и допустимы, поэтому входят в ответ.

Точки
\[
x=-3
\quad\text{и}\quad
x=0
\]
обращают знаменатель в ноль, поэтому не входят в ответ.

Точка
\[
x=7
\]
после сокращения исчезает из дроби, однако в исходном выражении
\[
x^2-7x=0,
\]
поэтому это значение обязательно исключается.

\resultlabel{ ошибка возникла не при выборе метода, а при раскрытии скобок после правильного приведения к общей дроби. Решение исходного неравенства:
\[
(-\infty;-4]\cup(-3;-1]\cup(0;7)\cup(7;+\infty).
\]
}

\textbf{Задача на закрепление.}

Решите неравенство
\[
\frac{2x+3-x^2}{x^2-4x}
\le
1+\frac{2}{x+2}.
\]

\vfill

\begin{center}
\small
Лёвин Артём Александрович эксклюзивно для Григория Архипова
\end{center}

% Краткое сообщение Григорию:
% Григорий, в работе ошибки встретились в четырёх местах.
% В задаче с векторами были найдены координаты нужного вектора,
% но не его длина. В задаче про ПЦР-тест были перепутаны вероятность
% положительного результата у заболевшего и вероятность заболевания
% при уже известном положительном результате. В задаче про теплоход
% уравнение было составлено правильно, но появилась арифметическая
% ошибка при дальнейших преобразованиях. В рациональном неравенстве
% верно выполнено приведение к общему знаменателю, но затем возникла
% ошибка при раскрытии скобок.

\end{document}